\section{2026 年 11 月核验增补}

本节按论文首次提交时间排列。每项方法附有可核验的外部入选依据。

% METHOD: 2410.07577
\paragraph{3D Vision-Language Gaussian Splatting：让语义使用独立的高斯合成权重}
2024.10~\cite{2410_07577}。颜色透明度用于图像成像，但由反射或半透明层解释的高斯对语义查询未必有用。该方法先在原语中融合视觉与语言特征，再通过语义专用的光栅化和经平滑的指标削弱无关原语。训练混合不同相机视图，减少语言特征依附单一相机或颜色的倾向。

原文在 LERF 评价定位和交并比，在 3D-OVS 比较开放词汇分割，并附有 Mip-NeRF 360 示例。语义参考来自 SAM 与 OpenCLIP，而不是完整三维语义标注，因此优势应按各个任务和相同特征协议解释。作者页列有 United Imaging Intelligence。

\textbf{入选依据} 论文署名单位包括 United Imaging Intelligence（\href{https://arxiv.org/html/2410.07577}{论文原文}）。

% METHOD: 2605.26503
\paragraph{不确定性感知高斯图：将语言导航的约束映射到三维价值}
2026.05~\cite{2605_26503}。该工作由全景 RGB-D 观察更新语义高斯地图，区分几何、语义与外观不确定性，并把这些信息融合为供导航策略使用的三维价值图，减少在证据薄弱位置作过强判断。

作者在 Matterport3D 模拟器的 R2R、RxR 和 REVERIE 评价导航与目标定位，并使用未见场景做模块消融。原文列出的局限包括静态仿真不覆盖移动对象、执行误差或实际传感器噪声，且结果不能外推至室外导航。作者单位包括 Huawei 2012 Labs。

\textbf{入选依据} 论文署名单位包括 Huawei（\href{https://arxiv.org/html/2605.26503v1}{论文原文}）。

% METHOD: 2606.07514
\paragraph{UniSHARP：在统一的射线距离空间预测单目高斯}
2026.06~\cite{2606_07514}。UniSHARP 将相机投影与高斯几何分开，用统一的射线距离表征不同成像模型，再以几何锚定高斯和特征条件残差描述所见场景。混合训练将透视、广角、鱼眼和全景输入纳入同一预测网络。

作者在多类相机数据用单源、多目标视图评价重建，并单独报告 WildRGB-D 上无需手工输入源相机内参的设置。无标定设置仍需目标相机定义评测视图；原文指出大幅改视点暴露大块不可见区域时可能出现空洞。作者单位包括 Insta360 Research。

\textbf{入选依据} 论文署名单位包括 Insta360（\href{https://arxiv.org/html/2606.07514v1}{论文原文}）。

% METHOD: 2607.04144
\paragraph{语义引导渐进移除：从跨视图参考补全高斯对象后的空洞}
2026.07~\cite{2607_04144}。该工作从多视图提取 DINOv2 特征，将语义相关的图像块用于修补对象遮挡过的区域，再按感知误差逐步细化质量较差的局部。最终编辑目标是高斯场中的物体移除与背景补全。

作者在前向场景和其他高斯修补方法间评价，并只对标记的移除区像素计算 PSNR、SSIM 与 LPIPS。这些指标衡量指定区域的生成图像质量，不能证明补全几何就是遮挡前的真实结构。作者单位包括字节跳动 PICO。

\textbf{入选依据} 论文署名单位包括 ByteDance（\href{https://arxiv.org/html/2607.04144v1}{论文原文}）。

% METHOD: 2609.11616
\paragraph{LangStreet：把语言证据绑定到锚点解码的持久地址}
2026.09~\cite{2609_11616}。LangStreet 将随视图生成的高斯原语用作路由，把多视图语义观察累加到持久槽位并向上汇成锚点摘要。对于观察稀疏的槽位，通过保持锚点方向的层级补全增强稳定性，再用相对锚点编码保存局部差异。

作者在 KITTI-360、Virtual KITTI 2 和 Waymo 上共享冻结的场景重建、解码器及二维语言引擎进行比较。论文显示紧凑字段可以接近完整字段并节省存储；结论对应同一预处理观察和持久结构，不能解释为重新训练重建或二维模型后的端到端提升。

\textbf{入选依据} 论文署名单位包括 vivo Mobile Communication Co., Ltd.（\href{https://arxiv.org/html/2609.11616v1}{论文原文}）。

% METHOD: 2609.18624
\paragraph{MoQSplat：将高斯场景按空间和质量层送入渐进传输}
2026.09~\cite{2609_18624}。MoQSplat 在 Media over QUIC 的内容组织中，将场景空间区域映射到 Track，再以语义分组和由粗到细的质量层组织对象。接收端按关注区域及层级决定请求优先级，发布者和转发者无需为各接收方维护单独的选层状态。

论文先测内容准备质量，再测网络条件变化下的端到端接收。两组使用的共同参照是 Mip-NeRF 360 的 Bicycle，内容准备以固定内参、距离和高度的绕行轨迹计算 VMAF。因此不透明度排序和带宽适配的观察属于单场景初步实验。

\textbf{入选依据} \href{https://github.com/emanuele-artioli/moqsplat}{公开实现}已找到与论文对应的代码文件。

% METHOD: 2609.28645
\paragraph{PePESeg3D：在几何和特征两阶段注入感知先验}
2026.09~\cite{2609_28645}。PePESeg3D 首先以单目深度和分割掩码引导高斯几何贴近结构边界，再通过深度颜色线索与跨视图质心约束学习多尺度语义，以补充二维掩码在部分粒度上的缺漏。

论文在 SPIn-NeRF、LERF-Mask 等设置报告标签传播与提示分割结果，另用 PSNR、SSIM 和 LPIPS 评价重建，两个任务的指标不能互代。作者观察到共享特征空间的尺度门控会使不同粒度的掩码互相影响，尤其在邻近尺度检索时。

\textbf{入选依据} \href{https://github.com/becow5x5/pepeseg3d}{公开实现}已找到与论文对应的代码文件。

% METHOD: 2609.34853
\paragraph{EviSplat：把多视图语义证据保留到查询时再聚合}
2026.09~\cite{2609_34853}。EviSplat 保留实例的各视图特征，并为高斯学习反映不同可见外观的紧凑证据场。查询时先按相关性选取实例观察，再用原语局部的可靠证据细化分数，不需要为了新查询重训表示。

论文在 LERF-OVS 和 3D-OVS 评价渲染掩码，在 ScanNet 按 OpenGaussian 设定评价 10 个场景的三维点，并将渲染后查语言特征与在三维先选高斯的方案分列。结论需区分这两类输出及其协议；保留逐视图证据也改变了语义字段的存储结构。

\textbf{入选依据} 论文署名单位包括 Huawei（\href{https://arxiv.org/abs/2609.34853}{论文原文}）。

% METHOD: 2610.08756
\paragraph{训练后语义提升：分别估计检测器、提升和表示误差}
2026.10~\cite{2610_08756}。该方法不重训既有高斯。给定单个目标类别，它累积各视图的目标和非目标证据，以高阈值选择种子，再以较低阈值接受含种子的连通区域。评估将预测映射到相机可见且有标注的网格顶点，从而排除不可见表面带来的混淆。

论文分别比较标注掩码和检测器掩码，在 Replica 选定设置并在 ScanNet++ 验证迁移。类若未出现在任何输入掩码就没有可提升的证据，且逐类处理可能让同一高斯落入多个类别，不能将该方法直接视为单次多类别分割。

\textbf{入选依据} \href{https://github.com/ivanver02/semantic-lifting-3dgs}{公开实现}已找到与论文对应的代码文件。

% METHOD: 2610.21619
\paragraph{目标感知高斯语言地图：最小特征注入与中心化定位}
2026.10~\cite{2610_21619}。该方法将指定目标在各视图中的视觉特征向文本嵌入靠拢，以闭式求出的最小强度注入，再将更新后的特征提升至高斯地图。查询时的中心化读出消除注入产生的共模分量，用以返回目标三维位置。

论文在三个基准将开放词汇分割与语言目标定位分开评价，并展示真实四足机器人。未做中心化的消融在三个数据上均可能因查询分数高度相关而退化；例如 ScanNet 的首选定位从基线 72.9 降至 30.0。论文的导航位置比较固定底层地图，不能等同于不同系统完整导航管线的比较。

\textbf{入选依据} 论文署名单位包括 Samsung（\href{https://arxiv.org/pdf/2610.21619v1}{论文原文}）。

% METHOD: 2611.19429
\paragraph{按查询分割高斯：从一个二维枢轴传播并投票到三维}
2026.11~\cite{2611_19429}。该方法以密集视觉语言特征给文本查询生成相似度图，选择二维枢轴区域，随后以视频分割器在有限视图预算内跟踪并细化掩码。根据每个高斯的可见证据投票，形成此次查询对应的三维实例掩码。

论文在 LERF-OVS 将首次查询所需时间、分割 mIoU 与定位准确率分别报告，并另列多个室内基准。该方案使用预训练模型和图像特征处理；视点变化大时跟踪可能漂移，严重遮挡也会使最终掩码不完整。

\textbf{入选依据} 论文署名单位包括 Qualcomm（\href{https://arxiv.org/abs/2611.19429}{论文原文}）。
